\scanpage{629}
\mspage{17}
En effet, en vertu de 2.1. $2^\circ$, cas B, \inserted{pour \deleted{\unsure{voir}} que $(W',C',F')$ soit un triple de Quillen clos,} il reste
à vérifier \deleted{que} (2.1.4\deleted{\unsure{b)}}\unsure{a)}). Mais cela résulte
de 1.8. appliqué aux deux couples de Quillen
spéciaux $(C,TF),(TC,F)$ dans $M$ : les couples
$(C',TC')$ et $(TC',F')$ sont des couples de Quillen
spéciaux, ce qui prouve en même temps
que le triple $(W',C',F')$ est spécial, \unsure{qed}.

Démontrons 2.1. \deleted{\unsure{Les}} les \inserted{premières deux} relations (2.1.3) sont
contenues dans 1.6, les relations $TC'\subset C'$, $TF'\subset F'$
sont tautologiques. Le critère (2.1.4) pour
que $(W',C',F')$ soit un triple de Quillen clos,
est contenu dans le rappel de la page 13, compte
tenu que l'on sait déjà que $TF'=F'\cap W'$
(relation (2.4 b)). Que la condition (1.1.4 b),
en présence de (2.1.4. a), soit équivalente
à la condition (2.1.5 b'), est contenu dans
la première assertion (2.5), \unsure{en notant|où notant} que
l'on a $W'=W'{}^{\natural}$ (ce qui résulte de la définition
de $W'$ et de $W^{\natural}=W^{\natural}$).

Il reste à prouver \inserted{$2^\circ)$, i.e} que la condition b) est
\deleted{\unsure{la forme équi}} satisfaite \deleted{\unsure{dès}} moyennant \inserted{a) et} b'') (2.1.6),
dans chacun des trois cas A), B), C).

\uline{Cas A} Par définition de $W'=(u^*)^{-1}(W)$, pour
prouver que \deleted{\unsure{u}} $TC'\subset W'$, il suffit de prouver

\[
\deleted{\unsure{(2.1.7)}}\qquad u^*(TC')\subset W'
\]

et a fortiori, il suffit de prouver

\[
u^*\deleted{\unsure{u}}(TC')\subset\deleted{\unsure{W}}TC\quad\text{i.e.}\quad u^*(TC')\longleftrightarrow F\tag*{(2.1.7)}
\]

Or par 1.5. appliqué au couple de foncteurs
% archive sheet number: 628
